\documentclass[12pt,a4paper]{article}
\usepackage[scheme=plain]{ctex}
\usepackage[margin=2.5cm]{geometry}
\usepackage{indentfirst}
\usepackage{setspace}
\usepackage{fontspec}
\setCJKmainfont[BoldFont=SimHei,ItalicFont=KaiTi]{SimSun}
\usepackage{newpxtext,newpxmath}

\usepackage{lastpage}
\usepackage{fancyhdr}
\pagestyle{fancy}
\renewcommand{\headrulewidth}{0pt}
\fancyhf{}
\fancyhead[C]{\small\itshape 阅读材料}
\fancyfoot[C]{\small \thepage}
\AtBeginDocument{\thispagestyle{fancy}}

\usepackage[nobottomtitles*]{titlesec}
\titleformat{\section}{\Large\bfseries}{\thesection}{1em}{}
\titlespacing*{\section}{0pt}{3ex}{2ex}

\newcommand{\clerknote}[1]{%
  \vspace{1.2em}%
  {\kaishu\small\textit{书记官按：#1}}%
}

\newcommand{\sourcenote}[1]{%
  \par\vspace{0.5em}%
  {\footnotesize\color{gray}（#1）}%
}

\usepackage{xcolor}
\usepackage{needspace}

\setlength{\parindent}{2em}
\setlength{\parskip}{0.4em}

\doublespacing

\title{\bfseries 试讲：准备！踏足于权力之野}
\author{}
\date{}

\begin{document}
\maketitle
\thispagestyle{fancy}
\vspace{-2em}

\begin{center}
\small\itshape 源自课程"重返城堡：如何在权力场域中缔结真实的人际关系？"
\end{center}

\vspace{1em}

% ========================
% 第一部分
% ========================
\section{来自 K 的第一封信}

同志们：

我到了。时间很少，说关键的。

城堡看不见。不是被什么挡住了，是它本身就不让你从远处辨认。但你往下走，走进村子，会发现这里每一件小事都指向它。老板看你时的眼神，人们沉默的方式，一个年轻人半夜把我叫醒盘问时那种既傲慢又紧张的样子——他在维护某种他自己也说不清楚的东西。还有一通夜里打到总办公厅的电话，有人接。这些就是我的全部线索。

我开始意识到一件事。这里最重要的东西不直观。它在语气里，在沉默的长度里，在人们视线交错的方式里。你得让自己完全沉浸进去感受，又随时能抽身出来看清楚自己感受到了什么。这件事比我预想的难得多。

出发前我们听到过不少传闻。有些调查队回来的人说，这个地方会压垮你，让你动弹不得，像被困在一张挣不脱的网上。也有人说，到这里的人会被自己的局限锁死，看什么都是隔着一层。传闻很多，众说纷纭。我没有到之前无法判断，到了以后也才几天，仍然不敢下结论。但有一点我可以告诉你们：到这里，用自己的眼睛看，在村子里扎下根，建立你们自己的关系——那时候你们看到的，跟传闻不会是一样的。

\begin{flushright}
K
\end{flushright}

\vspace{1em}
\newpage
{\large\bfseries 附件一}

卡夫卡在《审判》中写约瑟夫·K被押赴刑场时的心理："他想到苍蝇，它们那些细细的腿在挣脱粘纸时被折断了。"这种彻底的无能为力，这种面对不可理喻的境遇力量时的瘫痪感，渗透了他的全部作品。尽管《城堡》的情节走向与《审判》不同甚至相反，但这种从一只被困的、挣扎的苍蝇的视角看世界的感受，无处不在。这种体验，这种对被焦虑所主宰的世界的观感，这种人类任凭不可理解的恐怖摆布的处境，使卡夫卡的作品成为现代主义艺术的典型。那些在其他地方仅仅具有形式意义的技法，在这里被用来唤起面对一个全然陌生和敌意的现实时那种原始的敬畏。卡夫卡的焦虑，是现代主义最典型的体验。
\sourcenote{卢卡奇，《当代现实主义的意义》，1963，节选。由 GLM 5.2 据英文版译出。}

\vspace{1.5em}
{\large\bfseries 附件二}

\needspace{5\baselineskip}

K. 仿佛读过韦伯的理解社会学方法，在后者社会唯名论的立场上来看，社会制度并不是作为社会事实而存在，而是通过行动者对社会事实的解释而存在——K. 用来构建世界的数据恰是人们的解释和推测。显然，卡夫卡从未向我们展示任何关于城堡管理的确切知识，我们对城堡的唯一了解就是某些人的体验，通过外来者、下属和周边的人来了解城堡。这意味着，所有不可思议的、模棱两可的体验都是边缘人的真实体验。当然，如果我们将 K. 视为一位社会科学家，我们同样也要对其洞见的可靠性做出规范要求。K. 的"社会科学研究"存在两个问题：首先，他让自己被困住了，他以社会科学家的身份（以测量员的角色）进行实地研究，但再也没有离开过田野；其次，他显然局限于特定的人群和他们的经历。诚然，这并不完全是 K. 的错，他面对的可能是一个对被调查毫无兴趣的世界——事实上，许多社会科学研究者都经历过这种情况。在这样的背景下，我们不能排除研究者会通过改变观察视角与对象来发现秩序、系统与清晰的社会现象。
\sourcenote{约根森，"韦伯与卡夫卡：理性的与谜一般的官僚制", \textit{Public Administration}, 90: 194--210, 2012，节选。中译来自 Sociology 理论志。}

\clerknote{K 似乎向我们传递了某种目标和思想姿态。}

\newpage

% ========================
% 第二部分
% ========================
\section{如何做好田野}

记速记不是简单地在笔记本或电脑上写几个字的事。由于速记往往是在研究对象附近、甚至就在他们说话做事的当下写的，因此写速记本身就是一个社会性和互动性的过程。具体而言，研究者如何记、何时记速记，可能对那些人如何看待和理解她是谁、她在做什么，产生重要的影响。关于是否该记速记、如果记的话何时记怎么记，并没有什么硬性规则。但随着在某个场景中待的时间增长，通过不断试错的好处，田野研究者可以逐渐演化出一套独特的做法，使写速记这件事适应那个场景的轮廓和限制。

……

田野研究者还必须决定何时、何地、如何写速记。显然，低头看本子或键盘去写速记会分散研究者的注意力（即便只是一瞬间），使得对他人可能正在进行的复杂、快速而微妙的行动进行密切、持续的观察变得非常困难。但除了注意力的局限之外，记速记的决定还会对与田野中人们的关系产生重大影响。研究者努力与参与者建立紧密的联系，以便能被纳入他们生活中核心的活动。然而，在这些活动进行的过程当中，她可能体验到深深的矛盾：一方面，她可能希望通过趁话语刚说出口就记下来、趁场景正在上演就描摹下来，来保存当下的即时性；另一方面，她又可能觉得掏出一个笔记本或智能手机会毁掉这个瞬间，播下不信任的种子。参与者此时可能会把她看作一个主要兴趣在于发现他们的秘密、把他们最私密和最珍视的经历变成科学探究对象的人。

几乎所有民族志研究者在某些时刻都会感到被撕裂，一方面是自己的研究承诺，另一方面是自己真诚地融入所进入的那个世界的愿望。为了解决这些棘手的关系和道德问题，许多研究者主张：如果那些被研究的人没有完全知情和明确同意，就不应开展研究的任何环节。在这种观点看来，场景中的人必须被理解为合作者，他们与研究者积极合作，向外界讲述自己的生活和自己的文化。这种相互合作要求研究者请求许可才能撰写有关事件的内容，同时也要求尊重人们不愿揭示自己生活某些方面的意愿。

另一些田野研究者则觉得不必如此严格地被征求许可的要求所约束。一些人为这种立场辩护时坚持认为，田野研究者没有披露自己意图的特殊义务，因为所有的社会生活都包含掩饰的成分，没有人会完全暴露自己所有更深层的意图和私下的活动。另一些研究者指出，为自己而写的速记和田野笔记不会对他人造成直接伤害。这种做法当然是把那些艰难的道德和个人问题推迟到日后面对是否出版或以其他方式公开这些写作的决定时再去应对。最后，还有一些人主张向当地人隐瞒自己的研究目的，理由是所获得的信息将服务于更大的善。例如，如果研究者想要描述和公开无证工厂工人或养老院中老人的生存状况，他们就必须对控制这些场所准入权的当权者隐瞒自己的意图。
\sourcenote{Emerson, Fretz \& Shaw, \textit{Writing Ethnographic Fieldnotes}, 1995, 第二章。由 GLM 5.2 译。}

\clerknote{这本田野调查圣经为我们开展调查提供了策略上的建议，然而这些建议有待评估。}

\newpage

% ========================
% 第三部分
% ========================
\section{初入城堡}

{\small\sffamily 第一章，1--5页，K在城堡的第一夜}

K到达时，已经入夜了。村子被厚厚的积雪覆盖着。城堡所在的山岗连影子也不见，浓雾和黑暗包围着它，也没有丝毫光亮让人能约略猜出那巨大城堡的方位。K久久伫立在从大路通往村子的木桥上，举目凝视着眼前似乎是空荡荡的一片夜色。

然后他去找过夜的地方；酒店里人们还都没有睡，店老板虽然无房出租，但在对这位晚客的突然到来感到极度惊讶和惶乱之余，还是愿意让他在店堂里垫一个装稻草的口袋睡觉，K同意这一安排。有几个农民还在喝啤酒，但他无意同任何人交谈，便自己去阁楼上把草袋搬下来，在炉子附近躺下了。屋里很暖和，那几个农民说话声音很低，他用困倦的双眼打量了他们一番便倒头睡下了。

然而没有多久他便被吵醒了。这时只见一个城里人装束、长着一副演员般的面孔、浓眉细眼的年轻人同店老板一起站在他身边。那些农民也都还没有走，其中几个把椅子转过来对着他们，以便看得更清楚、听得更真切些。年轻人为吵醒了K而十分客气地向他道歉，自我介绍说，他是城堡主事的儿子，然后说道："这村子是城堡的产业，凡是在这里居住或过夜的，从某种意义上讲也算是在城堡居住或过夜，没有伯爵大人的许可，谁也不能在此居留。可是您并没有得到这样的许可，至少您并没有出示这样的证明嘛。"

K半坐起身子，捋了捋头发，仰头看着众人说道："我迷了路，这是摸到哪个村子来了？这里是有一座城堡吗？"

"那还用问，"年轻人慢条斯理地说，这时店堂各处都有人大惑不解地冲着K摇头。"这里是伯爵大人威斯特威斯的城堡。"

"一定要得到许可才能在这儿过夜吗？"K问道，似乎想弄清刚才他听到的那些话是不是在做梦。

答话是："是必须得到许可才行，"紧接着这个年轻人伸出胳臂，向店老板和酒客们问道："难道竟有什么人可以不必得到许可吗？"那话音和神态里，包含着对K的强烈嘲笑。

"那么我只好现在去讨要许可了。"K打着哈欠说，一面推开被子，似乎想站起来。

"向谁去讨要？"年轻人问。

"向伯爵大人，"K答道，"恐怕没有什么别的法子了吧。"

"现在，半夜三更去向伯爵大人讨要许可？"年轻人叫道，后退了一步。

"这不行吗？"K神色泰然地说，"那么您为什么叫醒我？"

这时年轻人憋不住火了。"真是活脱脱一副盲流口吻！"他喊叫起来，"伯爵衙门的尊严必须维护！我叫醒您是想告诉您：您必须立即离开伯爵领地！"

"好了，戏做够了吧。"K用异常轻的声音说，接着又躺下去，拉过被子盖在身上。"年轻人，您太过分了点，我明天还要再考虑考虑您今天的表现的。如果一定要见证的话，酒店老板和这里的各位先生就可以作证。现在请您听清楚：我是伯爵招聘来的土地测量员，明天我的助手就要带着仪器乘车随后跟来。我因为不想失去这个踏雪觅途的好机会，所以步行前来，可惜几次迷路，才到得这样晚。现在到城堡去报到时间已经太迟，这一点我自己很清楚，用不着您来赐教。正因为这样我才勉强在这草袋上凑合过夜，而您竟然——客气点说吧——举止失礼，打搅我休息。好了，我的话说完了，晚安，诸位先生！"说到这里K翻了一个身，转向炉子去了。

"土地测量员？"他听见背后将信将疑地发问，过后又无人做声了。然而不久那位年轻人便克制住自己，用压低了的——低到可以被看作是为了照顾K的睡眠，然而又大到能让他听清楚——的声音对老板说道："我现在就打电话去问一下。"怎么，在这个乡村小酒店里居然还有电话？唔，设备还真够齐全的。这些事，一件一件地听来也使K感到惊奇，不过总起来却又在他的意料之中。他发现，电话差不多正好摆在他的头顶上方，刚才他睡眼惺忪，没有注意到。现在，如果年轻人一定要打电话，那么他无论如何不能不打搅K的睡眠，问题只在于K让不让他打这个电话？K决定还是让他打。可是这样一来，在底下装睡便没有什么意思了，于是K又恢复了仰卧的姿势。他看见农民们怯怯地聚到一起，嘁嘁喳喳议论，看来土地测量员的到来不是一件小事。这时厨房的门早已打开，膀大腰圆的老板娘站在那里几乎把门框塞满，店老板踮着脚尖走到她跟前去告诉她刚才发生的一切。现在，电话接通，开始通话了，城堡主事已经就寝，是一位副主事——数位副主事当中某位名叫弗里茨的老爷——在那边接电话。年轻人自报姓名，说是叫施瓦尔策，接着便说他发现有一个名叫K的三十多岁的男子，衣冠不整，心安理得地在酒店里一个草袋上睡觉，枕着一个小得可怜的背囊，手边放着一根拐杖。他施瓦尔策自然觉得此人形迹可疑，而因为店主在这件事上显然失职，他施瓦尔策当然就责无旁贷地要过问此事，查明情况了。对于被叫醒盘问，对于他施瓦尔策按职责惯例作出的要将K逐出伯爵领地的警告，K的反应是很不耐烦，总之看起来火气不小，然而也许不无道理，因为他自称是伯爵大人聘来的土地测量员。在这种情况下，当然至少从例行手续上看有必要对他的话进行核实，因此，他施瓦尔策恳请弗里茨老爷在城堡总办公厅询问一下是否确有这样一位土地测量员应聘前来，并请将答复立即电话告知。

电话打完店堂里便安静下来，弗里茨在那边询问，这边在等着回话。K的神态一如既往，连头也不回，看样子一点也不急于知道结果如何，两眼茫然直视前方。施瓦尔策对事情的描述是一种不怀好意和谨小慎微的混合物，这番话给了K一种印象：城堡里连施瓦尔策这样的小人物也能很容易得到可说是外交手腕方面的训练。另外，看起来那里的人也决不偷懒；你看，总办公厅有人上夜班呢。而且显然很快就做出回答，因为这时弗里茨打电话来了。不过，好像这个回话太简短，因为施瓦尔策马上又气呼呼地把听筒挂上。"我早就说了嘛！"他大声叫道："土地测量员，连影子也没有！这人真是个卑鄙的、信口雌黄的流浪汉，说不定还更坏呢！"此时K脑子里闪过一个念头，他想：所有的人，就是施瓦尔策、那些农民，还有店老板、老板娘，眼看就要向他猛扑过来。为了至少避一避这个凶猛的势头，他从头到脚，整个儿蜷缩到被窝里面去了。这时电话铃又响起来，而且，K觉得声音特别急促响亮。于是他又把头从被子里慢慢伸出来。虽然电话几乎不可能仍是涉及K的事情，但所有的人还是一下子突然肃静下来，施瓦尔策则回到电话机旁去。他站在那儿耐心地听完了一番较长的解释之后，低声说道："那么是弄错了？我实在是太难为情。办公室主任亲自打来了电话？真是怪事，真是怪事。我该怎么向土地测量员先生解释才好呢？"

K听了这些话精神为之一振。这么说，城堡已经任命他为土地测量员了。这个情况一方面对他不利，因为它表明，城堡已经了解他的底细，在反复掂量了双方的力量对比之后，决定成竹在胸地开始同他较量。可是另一方面，情况又对他有利，因为在他看来这同时也证明对方低估了他，他将会有兴许比他自己起初敢于希冀的还要更多的一些自由。并且，如果以为通过这从心理战来看不能不说是高明的一着，即公开承认他的土地测量员身份，就能使他经常处于诚惶诚恐、心惊胆战的状态的话，那他们就错了；这一官方认可使他吃了一惊，但也就止此而已。

对羞怯地向他走来的施瓦尔策，K摆手示意他不必过来了；此刻人们忙不迭地恳求他赶快搬到老板的房里去住，他也拒绝了，只是接受了老板递给他的一杯催眠饮料，又从老板娘手里接过一盆水、肥皂和毛巾。现在根本用不着他开口叫人离开店堂，原来这时所有的人都把脸背过去使劲往外挤，可能是怕明天被他认出来吧。关灯之后，他终于可以休息了。他睡得很香，除了一两次被跑过的老鼠惊醒之外，一直睡到第二天早晨。

\clerknote{K 的第一晚是富有策略性的，但是 K 究竟做出了怎样的选择？}

\end{document}
